\documentclass{article}
\usepackage{amsmath}
\usepackage{amsfonts}
\usepackage{amssymb}
\usepackage[top=1.75cm,bottom=1.5cm,right=2cm,left=2cm]{geometry}
\setlength{\headheight}{22pt}
\usepackage{enumitem}
\usepackage{fancyhdr} % For customizing headers and footers

%---------------- Header
\pagestyle{fancy}
\fancyhf{} % Clear default header and footer
\fancyhead[L]{{\footnotesize CUNEF Universidad \\ Bachelor in Philosophy, Politics and Economics}} 
\fancyhead[R]{{\footnotesize A. G. Azze \\ \date{2025}}} 
\setlength{\headsep}{1.2cm}
\fancypagestyle{plain}{
  \fancyhf{}
  \fancyhead[L]{{\footnotesize CUNEF Universidad \\ Bachelor in Philosophy, Politics and Economics}} 
  \fancyhead[R]{{\footnotesize A. G. Azze \\ \date{2025}}} 
}

%---------------- Title
\title{Math Essentials \\ Unit 3 - Algebra \& Functions}
\author{}
\date{}

\begin{document}
\maketitle

\section*{Exercises}

\begin{enumerate}[label=\textbf{\arabic*.}, leftmargin=*, itemsep=0.6em]

\item A culture starts with $600$ bacteria and doubles every $3$ hours. Let $t$ be time in \emph{days}. Write a formula $N(t)$ for the population after $t$ days.

\item Simplify completely:
\[
\frac{t}{3t-1}+\frac{t}{3t+1}.
\]

\item Let $f(x)=3x+2$ and $g(x)=x^2-1$. Compute $g(f(x))$ and $f(g(x))$.

\item Let $f(x)=3x$. Decide whether $f(2a-3)-f(a-1)=f(a-2)$ holds for all $a\in\mathbb{R}$ and justify.

\item In right triangle $ABC$ with right angle at $C$, suppose $\cos(A)=\frac{4}{5}$. Find $\cos(B)$.

\item Apples cost is linear: $2$ lb cost $5.20$ euros and $5$ lb cost $13.00$ euros. Pears cost $C_P(p)=2.40\,p+0.50$ euros for $p$ pounds. Compute the per-pound cost of apples and pears, and state which is cheaper per pound.

\item Consider $h(x)=c\,(x-1)(x+4)^m$ with integers $m>0$ and $c\in\mathbb{Z}$. If $h(x)>0$ for $x\to+\infty$ and for $x\to-\infty$, what conditions must $m$ and $c$ satisfy? Explain briefly.

\item Expand and simplify: $(x-5)(x^2+4x-6)$.

\item Factor completely over $\mathbb{R}$: $4x^2-20x-24$.

\item Let $f(x)=a+bx$ be a line through $(3,8)$ and $(-1,4)$. Find $a$ and $b$.

\item Simplify $\displaystyle \frac{(1/x)}{y}$ and state whether it equals $\dfrac{y}{x}$. %for general nonzero $x,y$.

\item Build a quadratic function with zeros at $x=-2$ and $x=5$, and with second-order coefficient (this is the coefficient that multiplies $x^2$) equal to $-2$. Write $f(x)$ factorized and also give its expanded form.

\item Solve the inequality $|7-2x|\le 9$ and state the smallest $x$ that satisfies it.

\item Solve, for $x\in(-2,8)$: $|x-1|>3$.

\item If $\log_a x=b$ (with $a>0$, $a\ne1$, $x>0$), express $x$ in terms of $a$ and $b$.

\item Solve for $x$: $3^{\,2x-1}=20$.

\item 
\begin{enumerate}[label=(\alph*)]
\item If $(x+1)^2(y-5)>0$, what can you conclude about $y$? Are there any restrictions on $x$?
\item Solve for $x$: $16-x^2=(4-x)^2$.
\end{enumerate}

\item A theater has $n$ rows; each row has $5$ fewer seats than the number of rows. Write the total number of seats $T(n)$ as a function of $n$, and state the domain of $n$ for which this model makes sense.

\item If $x^{-1}y^{-1}=\dfrac14$, compute $x^3y^3$.

\item Simplify $\displaystyle \frac{x^{11}x^5}{x^8}$.

\item Find the value of $c$ for which the equation $4z^2+z=c$ has exactly one real solution.

\item Solve: $x^2-6x=7$.

\item State the range of $y=-3x^4+5$.

\item Solve for real $x$: $\sqrt{7x-2}-5=0$ (check the domain).

\item Simplify and state any restrictions:
\[
\frac{x+2y}{\,2y+x\,}+\frac{x-2y}{\,2y-x\,}.
\]

\item Create a linear equation whose \emph{only} solution is $x=6$.

\item Solve: $x^2-7x-5=0$.

\item Solve for $x$ (with $x\ne1$ and $x\ne -3$): $\displaystyle \frac{4}{x-1}=\frac{x}{x+3}$.

\end{enumerate}

\newpage
\section*{Solutions}

\begin{enumerate}[label=\textbf{\arabic*.}, leftmargin=*, itemsep=0.45em]

\item $N(t)=600\cdot 2^{(24/3)t}=600\cdot 2^{8t}$.

\item $\displaystyle \frac{t(3t+1)+t(3t-1)}{(3t-1)(3t+1)}=\frac{6t^2}{9t^2-1}$.

\item $g(f(x))=(3x+2)^2-1=9x^2+12x+3$; \quad $f(g(x))=3(x^2-1)+2=3x^2-1$.

\item $f(2a-3)-f(a-1)=3(2a-3)-3(a-1)=3a-6=f(a-2)$, so the identity holds for all $a$.

\item $A+B=90^\circ\Rightarrow \cos(B)=\sin (A)=\sqrt{1-(4/5)^2}=\frac{3}{5}$.

\item Apples: slope $=\dfrac{13-5.20}{5-2}=\dfrac{7.8}{3}=2.60$ €/lb. Pears: slope $=2.40$ €/lb. Pears are cheaper per pound.

\item Degree is $1+m$. Same sign at $\pm\infty$ requires even degree $\Rightarrow 1+m$ even $\Rightarrow m$ odd. To be positive at both ends, the leading coefficient must be $>0$, hence $c>0$.

\item $(x-5)(x^2+4x-6)=x^3- x^2-26x+30$.

\item $4x^2-20x-24=4(x^2-5x-6)=4(x-6)(x+1)$.

\item $b=\dfrac{8-4}{3-(-1)}=\dfrac{4}{4}=1$. Then $a+1\cdot 3=8\Rightarrow a=5$.

\item $\dfrac{(1/x)}{y}=\dfrac{1}{xy}$, which equals $\dfrac{y}{x}$ only iof $y = 1$ and $x \neq 0$; in general they are \emph{not} equal.

\item $f(x)=-2(x+2)(x-5)=-2x^2+6x+20$.

\item $|7-2x|\le 9\iff -9\le 7-2x\le 9 \iff -16\le -2x\le 2 \iff -1\le x\le 8$. Smallest $x=-1$.

\item $|x-1|>3\iff x< -2\ \text{or}\ x>4$. Intersect with $(-2,8)$ gives $(4,8)$.

\item $x=a^{\,b}$ (with $a>0$, $a\ne1$, $x>0$).

\item $(2x-1)\ln 3=\ln 20\ \Rightarrow\ x=\dfrac{1+\ln 20/\ln 3}{2}\approx 1.864$.

\item (a) Since $(x+1)^2\ge0$ and is $0$ only at $x=-1$, positivity forces $y-5>0$ and $x\ne -1$. Thus $y>5$ and $x\ne -1$.\\
(b) $16-x^2=(4-x)^2\Rightarrow 16-x^2=16-8x+x^2\Rightarrow 2x(x-4)=0\Rightarrow x=0$ or $x=4$.

\item $T(n)=n(n-5)$, meaningful for integers $n\ge 5$.

\item $(xy)^{-1}=\tfrac14\Rightarrow xy=4\Rightarrow x^3y^3=(xy)^3=64$.

\item $x^{11+5-8}=x^8$.

\item Discriminant $1-16c=0\Rightarrow c=\dfrac{1}{16}$.

\item $x^2-6x-7=0\Rightarrow x=\dfrac{6\pm\sqrt{36+28}}{2}=\dfrac{6\pm 8}{2}\Rightarrow x=7\ \text{or}\ x=-1$.

\item Since the quartic opens downward, $\max$ at $x=0$ is $5$. Range $(-\infty,5]$.

\item $\sqrt{7x-2}=5\Rightarrow 7x-2=25\Rightarrow x=\dfrac{27}{7}\approx 3.857$, which satisfies $x\ge \tfrac{2}{7}$.

\item $\dfrac{x+2y}{2y+x}=1$, and $\dfrac{x-2y}{2y-x}=\dfrac{x-2y}{-(x-2y)}=-1$, so the sum is $0$. Restrictions: $2y+x\ne0$ and $2y-x\ne0$.

\item $5x-30=0$ has the unique solution $x=6$.

\item $x=\dfrac{7\pm\sqrt{69}}{2}$.

\item $4(x+3)=x(x-1)\Rightarrow x^2-5x-12=0\Rightarrow x=\dfrac{5\pm\sqrt{73}}{2}$ (both admissible since $x\ne1,-3$).

\end{enumerate}

\end{document}
